\section*{Boolean payment of the root signed-KL rows}

For a Boolean function $f$ and a physical coordinate $i$, write
\[
a_A=\widehat f(A),\qquad b_A=\widehat f(A\cup\{i\}),
\qquad m_A=a_A^2+b_A^2,
\]
where $A$ avoids $i$.  Put
\[
C_A=\sum_{T\supseteq A}a_Tb_T,qquad
M_A=\sum_{T\supseteq A}m_T,qquad
\alpha_A={2a_Ab_A\over m_A},\qquad
\beta_A={2C_A\over M_A}.
\]
The exact restriction identity is
\[
\Delta_iU=B_i-\Gamma_i+K_i,qquad
K_i=\sum_Am_AD_2\left({1+\alpha_A\over2}\middle\|
                       {1+\beta_A\over2}\right).     \tag{1}
\]

\paragraph{Theorem.}
Let $K_i^{(0)}$ be the $A=\varnothing$ summand in (1).  Then
\[
\boxed{\sum_iK_i^{(0)}\le4\sum_i\widehat f(\{i\})^2
                         \le4I(f).}                  \tag{2}
\]
Consequently, if for some $a\ge4$ and $0\le b<1$ every tensor-prime
Boolean function obeys
\[
\sum_i\sum_{A\ne\varnothing}m_{i,A}D_2\left(
 {1+\alpha_{i,A}\over2}\middle\|{1+\beta_{i,A}\over2}\right)
\le(a-4)I+bU+\sum_i(\Gamma_i-B_i),                   \tag{3}
\]
then $G\le aI+bU$ for every Boolean function.

\paragraph{Proof.}
Boolean convolution at the nonzero shift $\{i\}$ gives
\[
C_\varnothing=\sum_Aa_Ab_A=0.                        \tag{4}
\]
Hence $\beta_\varnothing=0$.  Set
$a=\widehat f(\varnothing)$, $b=\widehat f(\{i\})$, and
$m=a^2+b^2$.  With the zero-mass convention,
\[
K_i^{(0)}=mD_2\left({1+2ab/m\over2}\middle\|{1\over2}\right)
=m\left[1-h_2\left({1+2ab/m\over2}\right)\right].   \tag{5}
\]
For $|x|\le1$,
\[
1-h_2\left({1+x\over2}\right)
={1\over\ln2}\sum_{k\ge1}{x^{2k}\over2k(2k-1)}
\le x^2.                                             \tag{6}
\]
Indeed all coefficients are positive and their sum is $\ln2$, as follows
by evaluating the series at $x=1$.  Applying (6) to $x=2ab/m$ yields
\[
K_i^{(0)}\le {4a^2b^2\over a^2+b^2}\le4b^2.         \tag{7}
\]
Summing (7) proves the first inequality in (2), and the second follows
because the singleton Fourier mass is part of total influence.

Summing (1), separating the root rows, and applying (2) gives
\[
G\le4I+\sum_i(B_i-\Gamma_i)+
 \sum_i\sum_{A\ne\varnothing}m_{i,A}D_2\left(
 {1+\alpha_{i,A}\over2}\middle\|{1+\beta_{i,A}\over2}\right).
\]
Thus (3) implies $G\le aI+bU$ on tensor-prime functions.  Exact tensor
additivity of $G,I,U$ then extends the same pair to every Boolean function.

\paragraph{Two-face compatibility.}
For distinct $i,j$, write on each Fourier two-face
\[
a_A=\widehat f(A),\ b_A=\widehat f(Ai),\
c_A=\widehat f(Aj),\ d_A=\widehat f(Aij).
\]
The convolution identities at shifts $\{i\},\{j\},\{i,j\}$ are
\[
\sum_A(a_Ab_A+c_Ad_A)=0,
\quad \sum_A(a_Ac_A+b_Ad_A)=0,
\quad \sum_A(a_Ad_A+b_Ac_A)=0.                      \tag{8}
\]
In particular the signed root correlation in each $j$-child for coordinate
$i$ is
\[
\sum_A(a_A\pm c_A)(b_A\pm d_A)=0,                   \tag{9}
\]
where (8) preserves both the ordinary and mixed terms.  The root payment is
therefore restriction-stable locally.  A global stopping argument must still
avoid repaying the same spectral degree mass at multiple nodes.

\paragraph{Scope.}
The theorem removes exactly the empty upper-set rows from the unpaid signed
KL aggregate.  It supplies no bound for $A\ne\varnothing$, no global
nonreuse theorem, and no universal pair $(a,b)$.  Universal GU and FEI
remain open.

